% abstract.tex

\chapter*{Abstract}

Zeroth-order optimization avoids gradients but incurs an additional dimension factor $d$ in the cost, making it impractical for high-dimensional ML. Computing function evaluations with $r$ parallel workers can reduce the wall-clock time by a factor of $r$, but existing theory relies on dense Gaussian perturbations requiring $O(dr)$ space to store. EGGROLL replaces these dense perturbations with structured low-rank perturbations $ab^\top$, dramatically reducing memory and computational cost, but prior work did not establish whether the resulting optimization procedure enjoys comparable convergence guarantees \citep{sarkar2025evolution}. We prove that, with high probability, rank-one EGGROLL converges to an $\ve$-first-order stationary point for $L$-smooth objectives in $\widetilde{O}\left(\frac{dm}{r}\cdot\frac{L(f(x_0)-f^*)}{\ve^2}\right)$ iterations. As a first step toward faster zeroth-order saddle escape, we analyze zeroth-order mirror descent. We then analyze a zeroth-order accelerated gradient descent algorithm yielding high-probability convergence to an $\ve$-optimal point in $\widetilde{O}\left(\frac{d}{r} \cdot \sqrt{\frac{L\Theta}{\ve}}\right)$ iterations. Collectively, these results show that low-rank perturbations are not merely an engineering trick: they retain high-probability guarantees while reducing memory requirements substantially.